% Options for packages loaded elsewhere
\PassOptionsToPackage{unicode}{hyperref}
\PassOptionsToPackage{hyphens}{url}
\PassOptionsToPackage{dvipsnames,svgnames,x11names}{xcolor}
%
\documentclass[
  11pt,
]{article}
\usepackage{amsmath,amssymb}
\usepackage{iftex}
\ifPDFTeX
  \usepackage[T1]{fontenc}
  \usepackage[utf8]{inputenc}
  \usepackage{textcomp} % provide euro and other symbols
\else % if luatex or xetex
  \usepackage{unicode-math} % this also loads fontspec
  \defaultfontfeatures{Scale=MatchLowercase}
  \defaultfontfeatures[\rmfamily]{Ligatures=TeX,Scale=1}
\fi
\usepackage{lmodern}
\ifPDFTeX\else
  % xetex/luatex font selection
\fi
% Use upquote if available, for straight quotes in verbatim environments
\IfFileExists{upquote.sty}{\usepackage{upquote}}{}
\IfFileExists{microtype.sty}{% use microtype if available
  \usepackage[]{microtype}
  \UseMicrotypeSet[protrusion]{basicmath} % disable protrusion for tt fonts
}{}
\makeatletter
\@ifundefined{KOMAClassName}{% if non-KOMA class
  \IfFileExists{parskip.sty}{%
    \usepackage{parskip}
  }{% else
    \setlength{\parindent}{0pt}
    \setlength{\parskip}{6pt plus 2pt minus 1pt}}
}{% if KOMA class
  \KOMAoptions{parskip=half}}
\makeatother
\usepackage{xcolor}
\usepackage[margin=2.2cm]{geometry}
\setlength{\emergencystretch}{3em} % prevent overfull lines
\providecommand{\tightlist}{%
  \setlength{\itemsep}{0pt}\setlength{\parskip}{0pt}}
\setcounter{secnumdepth}{5}
\usepackage{lmodern}
\usepackage[T1]{fontenc}
\usepackage{amsmath,amssymb}
\usepackage{newunicodechar}
\newunicodechar{}{\ensuremath{\bar\nu}}
\newunicodechar{}{\ensuremath{\hat\nu}}
\newunicodechar{}{\ensuremath{\tilde\sigma}}
\newunicodechar{}{\ensuremath{\bar v}}
\newunicodechar{}{\ensuremath{\hat\sigma}}
\newunicodechar{}{\ensuremath{\bar X}}
\newunicodechar{}{\ensuremath{\hat X}}
\newunicodechar{}{\ensuremath{\tilde W}}
\newunicodechar{}{\ensuremath{\dot\alpha}}
\newunicodechar{}{\ensuremath{\dot\beta}}
\newunicodechar{}{\ensuremath{\hat f}}
\newunicodechar{}{\ensuremath{\bar t}}
\newunicodechar{}{\ensuremath{\hat V}}
\newunicodechar{Ĉ}{\ensuremath{\hat C}}
\newunicodechar{ĉ}{\ensuremath{\hat c}}
\newunicodechar{ġ}{\ensuremath{\dot g}}
\newunicodechar{ᵀ}{\ensuremath{{}^{\mathsf{T}}}}
\newunicodechar{ℝ}{\ensuremath{\mathbb{R}}}
\newunicodechar{−}{\ensuremath{-}}
\newunicodechar{²}{\ensuremath{{}^{2}}}
\newunicodechar{³}{\ensuremath{{}^{3}}}
\newunicodechar{¹}{\ensuremath{{}^{1}}}
\newunicodechar{⁴}{\ensuremath{{}^{4}}}
\newunicodechar{⁸}{\ensuremath{{}^{8}}}
\newunicodechar{⁺}{\ensuremath{{}^{+}}}
\newunicodechar{⁻}{\ensuremath{{}^{-}}}
\newunicodechar{₂}{\ensuremath{{}_{2}}}
\newunicodechar{½}{\ensuremath{\tfrac12}}
\newunicodechar{¼}{\ensuremath{\tfrac14}}
\newunicodechar{·}{\ensuremath{\cdot}}
\newunicodechar{×}{\ensuremath{\times}}
\newunicodechar{÷}{\ensuremath{\div}}
\newunicodechar{±}{\ensuremath{\pm}}
\newunicodechar{σ}{\ensuremath{\sigma}}
\newunicodechar{φ}{\ensuremath{\varphi}}
\newunicodechar{ν}{\ensuremath{\nu}}
\newunicodechar{τ}{\ensuremath{\tau}}
\newunicodechar{π}{\ensuremath{\pi}}
\newunicodechar{ρ}{\ensuremath{\rho}}
\newunicodechar{λ}{\ensuremath{\lambda}}
\newunicodechar{μ}{\ensuremath{\mu}}
\newunicodechar{κ}{\ensuremath{\kappa}}
\newunicodechar{ω}{\ensuremath{\omega}}
\newunicodechar{β}{\ensuremath{\beta}}
\newunicodechar{α}{\ensuremath{\alpha}}
\newunicodechar{γ}{\ensuremath{\gamma}}
\newunicodechar{θ}{\ensuremath{\theta}}
\newunicodechar{η}{\ensuremath{\eta}}
\newunicodechar{χ}{\ensuremath{\chi}}
\newunicodechar{ψ}{\ensuremath{\psi}}
\newunicodechar{ξ}{\ensuremath{\xi}}
\newunicodechar{δ}{\ensuremath{\delta}}
\newunicodechar{ε}{\ensuremath{\varepsilon}}
\newunicodechar{Σ}{\ensuremath{\Sigma}}
\newunicodechar{Δ}{\ensuremath{\Delta}}
\newunicodechar{Λ}{\ensuremath{\Lambda}}
\newunicodechar{Π}{\ensuremath{\Pi}}
\newunicodechar{Ξ}{\ensuremath{\Xi}}
\newunicodechar{Γ}{\ensuremath{\Gamma}}
\newunicodechar{Φ}{\ensuremath{\Phi}}
\newunicodechar{Θ}{\ensuremath{\Theta}}
\newunicodechar{Ψ}{\ensuremath{\Psi}}
\newunicodechar{∫}{\ensuremath{\int}}
\newunicodechar{∂}{\ensuremath{\partial}}
\newunicodechar{√}{\ensuremath{\surd}}
\newunicodechar{∞}{\ensuremath{\infty}}
\newunicodechar{≈}{\ensuremath{\approx}}
\newunicodechar{≥}{\ensuremath{\ge}}
\newunicodechar{≤}{\ensuremath{\le}}
\newunicodechar{≠}{\ensuremath{\ne}}
\newunicodechar{≡}{\ensuremath{\equiv}}
\newunicodechar{→}{\ensuremath{\to}}
\newunicodechar{←}{\ensuremath{\leftarrow}}
\newunicodechar{⇒}{\ensuremath{\Rightarrow}}
\newunicodechar{⇔}{\ensuremath{\Leftrightarrow}}
\newunicodechar{↦}{\ensuremath{\mapsto}}
\newunicodechar{∈}{\ensuremath{\in}}
\newunicodechar{∝}{\ensuremath{\propto}}
\newunicodechar{∏}{\ensuremath{\prod}}
\newunicodechar{⊥}{\ensuremath{\perp}}
\newunicodechar{‖}{\ensuremath{\|}}
\newunicodechar{⟨}{\ensuremath{\langle}}
\newunicodechar{⟩}{\ensuremath{\rangle}}
\newunicodechar{∩}{\ensuremath{\cap}}
\newunicodechar{⊆}{\ensuremath{\subseteq}}
\newunicodechar{…}{\ensuremath{\ldots}}
\newunicodechar{̄}{}
\newunicodechar{̂}{}
\newunicodechar{̃}{}
\newunicodechar{̇}{}
\ifLuaTeX
  \usepackage{selnolig}  % disable illegal ligatures
\fi
\IfFileExists{bookmark.sty}{\usepackage{bookmark}}{\usepackage{hyperref}}
\IfFileExists{xurl.sty}{\usepackage{xurl}}{} % add URL line breaks if available
\urlstyle{same}
\hypersetup{
  pdftitle={Assignment Set 1 --- Defense Sheet},
  pdfauthor={Computational Finance (WI4151), TU Delft --- Prof.~Fang Fang},
  colorlinks=true,
  linkcolor={blue},
  filecolor={Maroon},
  citecolor={Blue},
  urlcolor={Blue},
  pdfcreator={LaTeX via pandoc}}

\title{Assignment Set 1 --- Defense Sheet}
\author{Computational Finance (WI4151), TU Delft --- Prof.~Fang Fang}
\date{}

\begin{document}
\maketitle

{
\hypersetup{linkcolor=blue}
\setcounter{tocdepth}{3}
\tableofcontents
}
\begin{quote}
Scope: this sheet covers the \textbf{coded} sub-questions only (2b, 2d,
2e, 3a, 3b, 3c). The hand-derivation parts (1a, 1b, 1c, 2a, 2c) are in
the PDF and are not repeated here, but Fang can jump from code to
derivation, so the cross-references below matter.
\end{quote}

\begin{center}\rule{0.5\linewidth}{0.5pt}\end{center}

\hypertarget{the-task-in-my-own-words}{%
\section{The task (in my own words)}\label{the-task-in-my-own-words}}

Assignment 1 has three blocks. The graded coded pieces sit in blocks 2
and 3.

\begin{itemize}
\tightlist
\item
  \textbf{Block 2 --- the COS method, sine variant.} The paper builds
  COS from a Fourier-\textbf{cosine} expansion; the assignment asks me
  to redo it with a Fourier-\textbf{sine} expansion.

  \begin{itemize}
  \tightlist
  \item
    \textbf{2b:} program the sine version of eq. (11) and reproduce the
    density-recovery errors of Table 1 of the COS paper for the standard
    normal density.
  \item
    \textbf{2d:} implement the sine-COS European \textbf{put} pricing
    formula (derived by hand in 2c) under Black--Scholes parameters
    \(S0=3, sigma=0.3, T=1, K=4, r=0.03\), and report the absolute error
    vs the analytic BS put as a function of \(N\).
  \item
    \textbf{2e:} extend the sine recovery one integration step further
    to recover the \textbf{CDF}, then verify it against the analytic
    log-normal CDF.
  \end{itemize}
\item
  \textbf{Block 3 --- importance sampling for a tail quantile.}
  Portfolio of \(N=1000\) digital options;
  \texttt{L\_T\ =\ sum\_i\ 1\{S\_i(T)\ \textgreater{}\ K\_i\}} with a
  common Brownian factor \(W\) and idiosyncratic factors \(B_i\).
  Estimate the 95\% quantile of \(L_T\).

  \begin{itemize}
  \tightlist
  \item
    \textbf{3a:} plain Monte Carlo.
  \item
    \textbf{3b:} importance sampling, mean shift \(W_T ~ N(1.5, T)\),
    reweight by \texttt{p/q}.
  \item
    \textbf{3c:} importance sampling, mean + variance shift
    \(W_T ~ N(1.5, 2T)\), reweight by \texttt{p/q}.
  \end{itemize}
\end{itemize}

\textbf{TODO (Bruno):} In one sentence each, state the \emph{point} of
each sub-question --- e.g.~why 2b uses the standard normal density
specifically, why 2e is ``one extra integration step'' beyond 2b/2c, and
what the 3b-vs-3c comparison is meant to demonstrate about importance
sampling.

\begin{center}\rule{0.5\linewidth}{0.5pt}\end{center}

\hypertarget{my-approach-and-why}{%
\section{My approach and why}\label{my-approach-and-why}}

What the code factually does (per block):

\begin{itemize}
\tightlist
\item
  \textbf{2b / 2e} build the sine series coefficients
  \texttt{F\_k\ =\ (2/(b-a))\ *\ Im\{\ phi(w\_k)\ exp(-i\ w\_k\ a)\ \}}
  with \texttt{w\_k\ =\ k*pi/(b-a)}, then sum \(F_k\) against either
  \texttt{sin(k\ pi\ (x-a)/(b-a))} (density, 2b) or the integrated sine
  \((1 - cos(...))\) form (CDF, 2e).
\item
  \textbf{2d} prices the put as
  \texttt{e\^{}\{-rT\}\ *\ sum\_k\ Im\{\ phi(w\_k)\ exp(-i\ w\_k\ a)\ \}\ *\ V\_k},
  where \(V_k\) is the analytic sine coefficient of the put payoff built
  from the \(phi_k\)/\(chi_k\) primitives.
\item
  \textbf{3a/3b/3c} simulate \(S_i(T)\) from the exact log-normal
  solution, count threshold crossings to form \(L_T\), and take a
  (weighted, in 3b/3c) empirical 95\% quantile.
\end{itemize}

\textbf{TODO (Bruno):} Justify, in your own words, the \emph{modelling}
choices that the code hard-codes: - 2b: why \([a,b] = [-10, 10]\) is a
safe truncation for the standard normal. - 2d: why you used the
\textbf{cumulant-based} range \([a,b] = c1 ± L*sqrt(c2 + sqrt(c4))\)
with \(L=10\), and why that is the ``correct'' way to set the range
rather than a fixed window. - 2e: why you switched to a \textbf{fixed}
window \([a,b] = [-2, 2]\) for the CDF test instead of reusing the
cumulant rule from 2d. - 3b/3c: why a mean shift of exactly \(1.5\), and
why doubling the variance is the natural ``second'' scheme to test.

\begin{center}\rule{0.5\linewidth}{0.5pt}\end{center}

\hypertarget{key-code-sections-file-what-it-does-why-i-wrote-it-that-way}{%
\section{Key code sections (file + what it does + why I wrote it that
way)}\label{key-code-sections-file-what-it-does-why-i-wrote-it-that-way}}

\hypertarget{b-2bcf-2.py-sine-density-recovery-standard-normal}{%
\subsection{\texorpdfstring{2b --- \texttt{2bCF\ (2).py} (sine density
recovery, standard
normal)}{2b --- 2bCF (2).py (sine density recovery, standard normal)}}\label{b-2bcf-2.py-sine-density-recovery-standard-normal}}

\begin{itemize}
\tightlist
\item
  \texttt{2bCF\ (2).py:4-5} set the truncation range \(a=-10, b=10\).
\item
  \texttt{2bCF\ (2).py:6-7} set the evaluation grid
  \(X_test = arange(-5,6)\) (i.e.~integers -5..5) and
  \(N_vals = [4,8,16,32,64]\).
\item
  \texttt{2bCF\ (2).py:10-12} \(normal_density(x)\) returns the standard
  normal pdf \texttt{(1/sqrt(2pi))\ exp(-x\^{}2/2)}.
\item
  \texttt{2bCF\ (2).py:14-16} \(char_function(w)\) returns
  \(exp(-0.5 w^2)\), the characteristic function of \(N(0,1)\).
\item
  \texttt{2bCF\ (2).py:18-33} \texttt{sin\_density(x\_vals,a,b,N)}:

  \begin{itemize}
  \tightlist
  \item
    \texttt{:22} \(k = arange(1, N+1)\) --- sine series starts at
    \(k=1\) (no \(k=0\) term).
  \item
    \texttt{:23} \texttt{wk\ =\ k*pi/(b-a)}.
  \item
    \texttt{:26}
    \texttt{fk\ =\ (2/(b-a))\ *\ Im\{\ char(wk)\ *\ exp(-i\ wk\ a)\ \}}
    --- the sine coefficients \(F_k\).
  \item
    \texttt{:30-31} for each \(x\), forms
    \texttt{sin(k\ pi\ (x-a)/(b-a))} and sums \(fk * sin_terms\).
  \end{itemize}
\item
  \texttt{2bCF\ (2).py:39-43} loops over \(N\), computes
  \texttt{max\textbar{}f\_approx\ -\ f\_true\textbar{}} and prints it.
\end{itemize}

\textbf{TODO (Bruno):} Explain \emph{why} line 26 takes the
\textbf{imaginary} part (tie it to the sine series and to Euler's
formula in your 2a derivation), and why the coefficient prefactor is
\texttt{2/(b-a)} and not \texttt{1/(b-a)}.

\hypertarget{d-2dcf-2.py-sine-cos-european-put-pricing}{%
\subsection{\texorpdfstring{2d --- \texttt{2dCF\ (2).py} (sine-COS
European put
pricing)}{2d --- 2dCF (2).py (sine-COS European put pricing)}}\label{d-2dcf-2.py-sine-cos-european-put-pricing}}

\begin{quote}
\textbf{FACT / FORMATTING FLAG:} this file is stored as a \textbf{single
physical line} (the newlines were stripped). So I cannot give honest
per-line numbers here --- I reference it by \textbf{function name}. Fang
may notice this; be ready to say it is a save/encoding artefact, not the
logic.
\end{quote}

\begin{itemize}
\tightlist
\item
  Parameters block: \(S0=3, sigma=0.3, T=1, K=4, r=0.03\);
  \texttt{x0\ =\ log(S0/K)}; \(mu = r - 0.5 sigma^2\).
\item
  Cumulant range: \(c1 = x0 + mu*T\), \(c2 = T*sigma^2\), \(c4 = 0\),
  \(L = 10\), \(a = c1 - L*sqrt(c2 + sqrt(c4))\),
  \(b = c1 + L*sqrt(c2 + sqrt(c4))\) (the comment cites ``eq 49'').
\item
  \(phi_BS(w, x, r, sigma, T)\): characteristic function of
  \texttt{y\ =\ ln(S\_T/K)}, returns
  \texttt{exp(i\ w\ mu\_y\ -\ 0.5\ var\_y\ w\^{}2)} with
  \(mu_y = x + (r - 0.5 sigma^2) T\), \(var_y = sigma^2 T\).
\item
  \(phi_k(k,a,b,c,d)\): returns
  \texttt{(1/wk)*(cos(wk(c-a))\ -\ cos(wk(d-a)))} --- the \(psi_k\)
  primitive from 2c.
\item
  \(chi_k(k,a,b,c,d)\): returns \(term(d)-term(c)\) with
  \texttt{term(y)\ =\ e\^{}y\ (sin(wk(y-a))\ -\ wk\ cos(wk(y-a)))\ /\ (1\ +\ wk\^{}2)}
  --- the \(chi_k\) primitive from 2c.
\item
  \texttt{V\_k\_put(a,b,K,ks)}: returns
  \texttt{(2K/(b-a))\ *\ (-chi\_k(ks,a,b,a,0)\ +\ phi\_k(ks,a,b,a,0))},
  i.e.~the payoff coefficient
  \texttt{V\_k\ =\ (2K/(b-a))(psi\_k(a,0)\ -\ chi\_k(a,0))}.
\item
  \texttt{sin\_put\_price(N,...)}: builds
  \texttt{cf\_part\ =\ Im\{\ phi\_BS(wk,\ x0,\ ...)\ exp(-i\ wk\ a)\ \}},
  multiplies by \(Vk\), and returns \(exp(-rT) * dot(cf_part, Vk)\).
\item
  \(bs_put(...)\): closed-form BS put for the reference price.
\item
  The driver loops \texttt{N\ in\ {[}4,8,16,32,64,128,256,512,1024{]}}
  and prints \(|SIN - BS|\).
\end{itemize}

\begin{quote}
\textbf{NAMING FLAG (be ready for this):} in the report (2c) \(psi_k\)
is the integral of \(sin\) and \(chi_k\) the integral of \(e^y sin\). In
this file the \textbf{function named \(phi_k\)} is the \(sin\) integral
(the report's \(psi_k\)) and the function named \(chi_k\) is the
\(e^y sin\) integral. So \(phi_k\) here = \(psi_k\) in the report. The
numbers still match because \texttt{V\_k\_put} pairs them correctly, but
the name \(phi_k\) collides with the characteristic function \(phi\).
\end{quote}

\textbf{TODO (Bruno):} Justify the \emph{content} choices: why
\(c4 = 0\) is acceptable here (what is the 4th cumulant of a Gaussian
log-price?), why \texttt{x0\ =\ log(S0/K)} is the right state variable,
and why the discount factor is \(e^{-rT}\) rather than
\(e^{-r*Delta t}\) with some other \(Delta t\).

\hypertarget{e-2ecf-2.py-sine-cos-cdf-recovery}{%
\subsection{\texorpdfstring{2e --- \texttt{2eCF\ (2).py} (sine-COS CDF
recovery)}{2e --- 2eCF (2).py (sine-COS CDF recovery)}}\label{e-2ecf-2.py-sine-cos-cdf-recovery}}

\begin{itemize}
\tightlist
\item
  \texttt{2eCF\ (2).py:5-8} parameters \(r=0.03, S0=3, sigma=0.3, T=1\).
\item
  \texttt{2eCF\ (2).py:11-12} \(mu = (r - 0.5 sigma^2) T\),
  \(sigma_aux = sigma*sqrt(T)\) --- moments of \texttt{ln(S\_T/S\_0)}.
\item
  \texttt{2eCF\ (2).py:14-15} fixed range \(a=-2, b=2\).
\item
  \texttt{2eCF\ (2).py:16-17} \(X_test = linspace(-0.5, 0.5, 11)\),
  \(N_vals = [4,8,16,32,64,128]\).
\item
  \texttt{2eCF\ (2).py:19-23} \(char_function(w)\) = CF of
  \texttt{ln(S\_T/S\_0)}:
  \texttt{exp(i\ w\ mu\ -\ 0.5\ sigma\_aux\^{}2\ w\^{}2)}.
\item
  \texttt{2eCF\ (2).py:25-29}
  \texttt{F\_exact(z)\ =\ norm.cdf((z\ -\ mu)/sigma\_aux)} --- analytic
  normal CDF.
\item
  \texttt{2eCF\ (2).py:31-46} \texttt{sin\_cdf(z\_vals,N,a,b)}:

  \begin{itemize}
  \tightlist
  \item
    \texttt{:36-37} \(ks = arange(1,N+1)\),
    \texttt{wk\ =\ ks\ pi/(b-a)}.
  \item
    \texttt{:40} \(im = Im{ char(wk) exp(-i wk a) }\).
  \item
    \texttt{:43}
    \texttt{cosine\ =\ cos(\ outer((z-a)/(b-a),\ ks)\ *\ pi\ )} ---
    matrix of \texttt{cos(k\ pi\ (z-a)/(b-a))}.
  \item
    \texttt{:46} returns
    \texttt{(2/pi)\ *\ (im/ks)\ @\ (1\ -\ cosine).T}, i.e.
    \texttt{F(z)\ \textasciitilde{}=\ (2/pi)\ sum\_k\ (1/k)\ Im\_k\ (1\ -\ cos(k\ pi\ (z-a)/(b-a)))}.
  \end{itemize}
\item
  \texttt{2eCF\ (2).py:54-57} loop over \(N\), print
  \texttt{max\textbar{}F\_sin\ -\ F\_true\textbar{}}.
\end{itemize}

\textbf{TODO (Bruno):} Explain \emph{why} the CDF formula carries a
\texttt{1/k} factor and a \texttt{2/pi} prefactor (where do they come
from when you integrate the density term by term?), and why the
\((1 - cos(...))\) shape is the integrated form of \(sin(...)\).

\hypertarget{a-3acf-2.py-plain-monte-carlo-quantile}{%
\subsection{\texorpdfstring{3a --- \texttt{3aCF\ (2).py} (plain Monte
Carlo
quantile)}{3a --- 3aCF (2).py (plain Monte Carlo quantile)}}\label{a-3acf-2.py-plain-monte-carlo-quantile}}

\begin{itemize}
\tightlist
\item
  \texttt{3aCF\ (2).py:5-12} parameters
  \texttt{N\_stocks=1000,\ S\_0=1,\ r=0.01,\ sigma1=0.8,\ sigma2=0.6,\ T=1,\ K=2},
  \(alpha = 1-0.95\).
\item
  \texttt{3aCF\ (2).py:14}
  \texttt{drift\ =\ r\ -\ 0.5*(sigma1\^{}2\ +\ sigma2\^{}2)}.
\item
  \texttt{3aCF\ (2).py:16-31} \(simulate_plainMC(paths)\):

  \begin{itemize}
  \tightlist
  \item
    \texttt{:22} common factor \(Wt ~ N(0, sqrt(T))\) (shape \(paths\)).
  \item
    \texttt{:23} idiosyncratic \(Bi ~ N(0, sqrt(T))\) (shape
    \(paths x N_stocks\)).
  \item
    \texttt{:26}
    \texttt{S\_T\ =\ S\_0\ exp(drift*T\ +\ sigma1*Wt{[}:,None{]}\ +\ sigma2*Bi)}.
  \item
    \texttt{:29} \texttt{L\_T\ =\ sum(S\_T\ \textgreater{}\ K,\ axis=1)}
    --- count of in-the-money names per path.
  \end{itemize}
\item
  \texttt{3aCF\ (2).py:36} \texttt{np.random.seed(42)}.
\item
  \texttt{3aCF\ (2).py:40-44} loops \(paths\) over \([1000..500000]\),
  \texttt{q\ =\ np.quantile(L\_T,\ 1-alpha)} = 0.95 quantile.
\item
  \texttt{3aCF\ (2).py:47-55} semilog convergence plot.
\end{itemize}

\textbf{TODO (Bruno):} Justify why
\texttt{np.std}/\texttt{np.random.normal(0,\ sqrt(T))} (a \(N(0,T)\)
draw) is the \emph{exact} terminal law of \(W_T\) here, why no
time-stepping is needed, and why the empirical quantile is taken at
\(1 - alpha = 0.95\) given the assignment's definition
\texttt{P(L\_T\ \textgreater{}=\ q\_alpha)\ \textless{}=\ alpha}.

\hypertarget{b-3bcf-2.py-importance-sampling-mean-shift}{%
\subsection{\texorpdfstring{3b --- \texttt{3bCF\ (2).py} (importance
sampling, mean
shift)}{3b --- 3bCF (2).py (importance sampling, mean shift)}}\label{b-3bcf-2.py-importance-sampling-mean-shift}}

\begin{itemize}
\tightlist
\item
  \texttt{3bCF\ (2).py:17-18} IS parameters \(mu_shift = 1.5\),
  \texttt{mu\_shiftsq\ =\ mu\_shift\^{}2\ /\ (2T)}.
\item
  \texttt{3bCF\ (2).py:20-34} \(weighted_quantile(values, weights, q)\):

  \begin{itemize}
  \tightlist
  \item
    \texttt{:25-27} sort \(values\), reorder \(weights\) to match.
  \item
    \texttt{:30} \texttt{cdf\ =\ cumsum(w\_sorted)/sum(w\_sorted)}.
  \item
    \texttt{:33-34} \(idx = searchsorted(cdf, q)\), return
    \(v_sorted[idx]\).
  \end{itemize}
\item
  \texttt{3bCF\ (2).py:43}
  \texttt{W\_q\ \textasciitilde{}\ N(mu\_shift,\ sqrt(T))} --- common
  factor drawn under \(q\).
\item
  \texttt{3bCF\ (2).py:44} \(Bi ~ N(0, sqrt(T))\) --- idiosyncratic
  factors unchanged.
\item
  \texttt{3bCF\ (2).py:47}
  \texttt{S\_T\ =\ exp(drift*T\ +\ sigma1*W\_q{[}:,None{]}\ +\ sigma2*Bi)}
  (note: no explicit \(S_0\), since \(S_0=1\)).
\item
  \texttt{3bCF\ (2).py:50}
  \texttt{L\_T\ =\ sum(S\_T\ \textgreater{}\ K,\ axis=1)}.
\item
  \texttt{3bCF\ (2).py:53}
  \texttt{ratio\ =\ exp(-(mu\_shift/T)*W\_q\ +\ mu\_shiftsq)} ---
  likelihood ratio \texttt{p/q}.
\item
  \texttt{3bCF\ (2).py:56}
  \texttt{q\_95\ =\ weighted\_quantile(L\_T,\ ratio,\ 0.95)}.
\end{itemize}

\textbf{TODO (Bruno):} Derive the \(ratio\) on line 53 from \texttt{p/q}
of two Gaussians and confirm it equals the report's
\(exp(-1.5 w + 1.125)\) at \(T=1\). Also justify \emph{why only \(W\) is
reweighted and not \(B_i\)} (the report says the \(B_i\) ratio is 1 ---
state precisely why).

\hypertarget{c-3ccf-2.py-importance-sampling-mean-variance-shift}{%
\subsection{\texorpdfstring{3c --- \texttt{3cCF\ (2).py} (importance
sampling, mean + variance
shift)}{3c --- 3cCF (2).py (importance sampling, mean + variance shift)}}\label{c-3ccf-2.py-importance-sampling-mean-variance-shift}}

\begin{itemize}
\tightlist
\item
  \texttt{3cCF\ (2).py:17-18} IS parameters \(mu_shift = 1.5\),
  \(var_q = 2*T\).
\item
  \texttt{3cCF\ (2).py:20-34} \(weighted_quantile\) --- identical
  structure to 3b.
\item
  \texttt{3cCF\ (2).py:43}
  \texttt{W\_q\ \textasciitilde{}\ N(mu\_shift,\ sqrt(var\_q))} =
  \(N(1.5, 2T)\).
\item
  \texttt{3cCF\ (2).py:44,47,50} same \(Bi\), \(S_T\), \(L_T\)
  construction as 3b.
\item
  \texttt{3cCF\ (2).py:54}
  \texttt{ratio\ =\ sqrt(2)*exp((-W\_q\^{}2\ -\ 3\ W\_q\ +\ 2.25)/(4T))}
  --- the \texttt{p/q} for the mean+var shift.
\item
  \texttt{3cCF\ (2).py:56}
  \texttt{q\_95\ =\ weighted\_quantile(L\_T,\ ratio,\ 0.95)}.
\end{itemize}

\textbf{TODO (Bruno):} Derive the line-54 ratio from \(p ~ N(0,T)\),
\(q ~ N(1.5, 2T)\) and confirm the \texttt{sqrt(2)} prefactor comes from
the variance mismatch in the normalising constants. Then explain why
this scheme is \emph{worse} for the 95\% quantile than 3b (the report's
argument about mass landing too far in the tail) --- restate it as your
own.

\begin{center}\rule{0.5\linewidth}{0.5pt}\end{center}

\hypertarget{design-decisions-i-must-justify}{%
\section{Design decisions I must
justify}\label{design-decisions-i-must-justify}}

\begin{itemize}
\tightlist
\item
  \textbf{Sine vs cosine expansion.} The code implements the sine series
  throughout (coefficients via \(Im{...}\)). \textbf{TODO (Bruno):} why
  does the assignment want sine, and what changes structurally vs the
  cosine COS (the \(k=0\) half-weight term, the \(Im\) vs \(Re\))?
\item
  \textbf{Truncation range, three different choices.} 2b uses
  \([-10,10]\), 2d uses the cumulant rule with \(L=10\), 2e uses
  \([-2,2]\). \textbf{TODO (Bruno):} justify each range and explain
  whether \([-2,2]\) in 2e is wide enough for \(N(mu, sigma_aux^2)\)
  with \(mu = -0.015\), \(sigma_aux = 0.3\) (it is \textasciitilde±6.7
  sigma --- confirm and say so). This is exactly the kind of ``is your
  range safe?'' question Fang asks.
\item
  \textbf{\(L = 10\) cumulant multiplier.} \textbf{TODO (Bruno):} the
  paper uses \(L\) around 8--10; justify your \(L=10\) and what error
  you'd see if \(L\) were too small (link to your Assignment 2
  short-maturity / small-c2 bug).
\item
  \textbf{\texttt{np.quantile} interpolation vs the
  \(weighted_quantile\) searchsorted rule.} 3a uses NumPy's
  \texttt{np.quantile} (linear interpolation); 3b/3c use a custom
  step-function \(searchsorted\) quantile. \textbf{TODO (Bruno):} why
  two different quantile estimators, and does the mismatch bias the
  3a-vs-3b/3c comparison?
\item
  \textbf{\(seed(42)\).} Used in 3a/3b/3c for reproducibility.
  \textbf{TODO (Bruno):} confirm this is only for reproducibility and
  does not, e.g., reuse the same \(W\) across schemes in a way that
  matters.
\end{itemize}

\begin{center}\rule{0.5\linewidth}{0.5pt}\end{center}

\hypertarget{results-and-what-they-mean}{%
\section{Results and what they mean}\label{results-and-what-they-mean}}

Reported in the PDF (factual, from the tables):

\begin{itemize}
\tightlist
\item
  \textbf{2b:} max density error falls
  \(2.1e-1 (N=4) -> 1.0e-16 (N=64)\) --- matches the paper's Table 1
  order of magnitude; spectral convergence.
\item
  \textbf{2d:} \(|SIN - BS|\) falls to \(~3.3e-16\) by \(N=64\) and
  stays there; BS reference put \(= 0.9918191782\).
\item
  \textbf{2e:} max CDF error falls \(1.37e-1 (N=4) -> 3.67e-11 (N=32)\)
  and then \textbf{plateaus} at \(3.67e-11\).
\item
  \textbf{3a:} 95\% quantile estimate settles near \(590\) as \(paths\)
  grows to \(5e5\).
\item
  \textbf{3b:} estimate converges to \(~587\) and visibly stabilises
  with far fewer paths than 3a.
\item
  \textbf{3c:} estimate ends near \(589\) but the path is noisier; the
  report argues it is \emph{less} efficient than 3b for this quantile.
\end{itemize}

\textbf{TODO (Bruno):} - 2d/2e: explain the \textbf{plateau} (why does
the error stop improving --- round-off vs truncation-range error?). For
2e specifically, why does it plateau at \(3.67e-11\) and not at machine
\(1e-16\) like 2d? This is the single most likely ``your convergence
stalled --- why?'' question. - 3a/3b/3c: are 590 / 587 / 589
\emph{consistent} (do their confidence intervals overlap)? What is the
``true'' quantile and how confident are you?

\begin{center}\rule{0.5\linewidth}{0.5pt}\end{center}

\hypertarget{the-hardest-question-fang-could-ask-here-and-my-answer}{%
\section{The hardest question Fang could ask here --- and my
answer}\label{the-hardest-question-fang-could-ask-here-and-my-answer}}

\textbf{Likely hardest:} \emph{``Your 2e CDF error stalls at 3.67e-11
while your 2d price reaches 1e-16. Both use the sine-COS machinery. Why
the difference, and is 3.67e-11 a bug or expected?''}

This bites because it sits exactly on the truncation-range / convergence
theme Fang co-authored, and the two tables disagree by five orders of
magnitude.

\textbf{TODO (Bruno):} Write the answer. Components you must cover: 1.
The CDF formula integrates the density term by term, so its error
inherits the \textbf{truncation-range} error \(f(a), f(b) != 0\) at the
endpoints --- i.e.~the residual comes from \([a,b]=[-2,2]\) not being
infinitely wide, not from too few terms \(N\). (Confirm by saying what
happens if you widen \([a,b]\).) 2. Why 2d reaches machine precision
instead: the put payoff coefficients \(V_k\) decay and the BS density is
extremely well captured by the cumulant range, so series-truncation
error drops below round-off. 3. State plainly whether \(3.67e-11\) is
acceptable or whether widening \([-2,2]\) would push it lower.

(Second candidate hardest: \emph{``Derive line 53 / line 54 likelihood
ratios from scratch and tell me why the variance shift in 3c hurts.''}
--- see the 3b/3c TODOs.)

\begin{center}\rule{0.5\linewidth}{0.5pt}\end{center}

\hypertarget{show-me-in-your-code-where-you-do-x-anticipated-spots}{%
\section{``Show me in your code where you do X'' --- anticipated
spots}\label{show-me-in-your-code-where-you-do-x-anticipated-spots}}

\begin{itemize}
\tightlist
\item
  \textbf{``\ldots where you build the sine coefficients \(F_k\).''}
  -\textgreater{} \texttt{2bCF\ (2).py:26}; same construction at
  \texttt{2eCF\ (2).py:40} and inside \texttt{sin\_put\_price}
  (\(cf_part\)) in \texttt{2dCF\ (2).py}.
\item
  \textbf{``\ldots where you set the COS truncation range.''}
  -\textgreater{} \texttt{2bCF\ (2).py:4-5} (fixed), cumulant block in
  \texttt{2dCF\ (2).py} (\(a = c1 - L*sqrt(c2+sqrt(c4))\)),
  \texttt{2eCF\ (2).py:14-15} (fixed \([-2,2]\)).
\item
  \textbf{``\ldots where the characteristic function enters pricing.''}
  -\textgreater{} \(phi_BS\) and \(cf_part\) in \texttt{2dCF\ (2).py}.
\item
  \textbf{``\ldots where you compute the payoff coefficients \(V_k\).''}
  -\textgreater{} \texttt{V\_k\_put} (with \(phi_k\)/\(chi_k\)) in
  \texttt{2dCF\ (2).py}.
\item
  \textbf{``\ldots where you integrate the sine to get the CDF.''}
  -\textgreater{} \texttt{2eCF\ (2).py:43-46} (the \((1 - cosine)\)
  term).
\item
  \textbf{``\ldots where you get the exact terminal stock price.''}
  -\textgreater{} \texttt{3aCF\ (2).py:26}, \texttt{3bCF\ (2).py:47},
  \texttt{3cCF\ (2).py:47}.
\item
  \textbf{``\ldots where you form \(L_T\).''} -\textgreater{}
  \texttt{3aCF\ (2).py:29}, \texttt{3bCF\ (2).py:50},
  \texttt{3cCF\ (2).py:50}.
\item
  \textbf{``\ldots where you apply the likelihood ratio.''}
  -\textgreater{} \texttt{3bCF\ (2).py:53}, \texttt{3cCF\ (2).py:54}.
\item
  \textbf{``\ldots where you take the weighted quantile.''}
  -\textgreater{} \(weighted_quantile\) in \texttt{3bCF\ (2).py:20-34} /
  \texttt{3cCF\ (2).py:20-34}, called at \texttt{:56}.
\item
  \textbf{``\ldots where you take the plain quantile.''} -\textgreater{}
  \texttt{3aCF\ (2).py:42} (\texttt{np.quantile(L\_T,\ 1-alpha)}).
\end{itemize}

\end{document}
